\documentclass[10pt,a4paper]{amsart}
\usepackage[margin=19mm]{geometry}
\usepackage{amsmath,amssymb,mathtools}
\usepackage[scr=rsfs,cal=euler]{mathalfa}
\usepackage{xfrac}
\usepackage[unicode,hidelinks,bookmarks=false]{hyperref}
\newcommand{\ie}{i.e.\ }
\newcommand{\cf}{cf.\ }
\newcommand{\resp}{respectively\ }
\newcommand{\Subsection}{\subsection}
\providecommand{\nobreakdash}{\nobreak\hbox{-}}
\setlength{\emergencystretch}{2em}
\multlinegap0pt
\usepackage{amssymb}
\usepackage{tikz}
\newtheorem{lem}{Lemma}
\title{V-numbers of symbolic powers of cover ideals of graphs}
\author{Thanh Vu}
\date{Source: arXiv:2608.15268v1 (CC BY 4.0)}
\begin{document}
\maketitle

\noindent\emph{Source and licence.} V-numbers of symbolic powers of cover ideals of graphs. Version arXiv:2608.15268v1 (CC BY 4.0), distributed by arXiv under the Creative Commons Attribution 4.0 International licence, http://creativecommons.org/licenses/by/4.0/, as recorded in the arXiv licence metadata for this exact version and shown on the abstract page https://arxiv.org/abs/2608.15268v1. Full licence notice: \texttt{licenses/arxiv-2608.15268-CC-BY-4.0.txt}.\par\smallskip

\noindent\textit{Editorial scope.} Counted unit: the article's lem lem\_decomp1 (source0.tex lines 292--294). The article's complete original proof of it is delivered with it (source0.tex lines 296--323, one \texttt{proof} environment of 28 lines). No mathematics is written, repaired or re-derived: the statement and the proof are the article's own lines, in the article's own notation, and no retained line is substituted or re-flowed.  No further source-local context is needed: every cross-reference of every retained line resolves inside the two delivered spans.  Nothing was omitted from any retained span, so the package carries no omission record.  No retained line contains a citation, so the wrapper carries no bibliography and no bibliography entry.  No retained line contains a figure environment, an image-inclusion command or drawing code, so no image is delivered and the package declares no asset dependency.  Editorial formatting only, in addition: an AMS \texttt{amsart} preamble that restates the article's own declarations and macros used by the retained lines, namely the environments \texttt{lem} on separate counters, together with the standard packages those lines need.  The retained environments print in this document as Lemma 1 in order of appearance, whereas the article prints the same items with the numbering of its own full document.  The complete native source of the article as distributed in the arXiv e-print for this exact version is bundled unchanged as \texttt{sources/207726/source0.tex}.
\begin{lem}\label{lem_decomp1} 
Let $G$ be a bipartite graph, and let $P = \{i,j\}$ be an edge of $G$. Assume that $t \ge 1$ is an integer. Let $\mathbf{c} \in \mathbb{N}^n$ be a $P$-almost $(t+1)$-cover of $G$. Then there exists a $1$-cover $\mathbf{b}$ and a $P$-almost $t$-cover $\mathbf{a}$ of $G$ such that $\mathbf{c} = \mathbf{a} + \mathbf{b}$.
\end{lem}

\begin{proof}
Let $V(G) = U \cup V$ be a bipartition of $G$, and assume without loss of generality that $i \in U$ and $j \in V$. Choose $r \in \{1, \ldots, t+1\} \setminus \{c_i + 1\}$. Define $\mathbf{b}$ on $U$ and $V$ as follows:
\[
b_u = \begin{cases} 
1 & \text{if } u \in U \text{ and } c_u \ge r, \\
0 & \text{otherwise,} 
\end{cases}
\]
and
\[
b_v = \begin{cases} 
1 & \text{if } v \in V \text{ and } c_v \ge t + 2 - r, \\
0 & \text{otherwise.} 
\end{cases}
\]
We now prove that $\mathbf{b}$ is a $1$-cover of $G$.

Indeed, suppose for the sake of contradiction that there exists an edge $\{u,v\} \in E(G)$ with $u \in U$ and $v \in V$ such that $b_u = 0$ and $b_v = 0$. By definition, we have $c_u \le r - 1$ and $c_v \le t + 1 - r$. Consequently, $c_u + c_v \le t$. Since $\mathbf{c}$ is a $P$-almost $(t+1)$-cover, $c_u + c_v \ge t+1$ for all edges $\{u,v\} \neq \{i,j\}$. Thus, $c_u + c_v \le t$ is possible if and only if $u = i$ and $v = j$. But then $c_i \le r - 1$ and $c_j \le t + 1 - r$. Combined with $c_i + c_j = t$, this forces $c_i = r - 1$, which contradicts our choice of $r \neq c_i + 1$. Hence, no such edge exists, and $\mathbf{b}$ is a $1$-cover of $G$.

Now, let $\mathbf{a} = \mathbf{c} - \mathbf{b}$. We show that $\mathbf{a}$ is a $P$-almost $t$-cover of $G$. From the construction, exactly one of $b_i, b_j$ is equal to $1$, so $b_i + b_j = 1$. Thus, $a_i + a_j = (c_i + c_j) - (b_i + b_j) = t - 1$. 

For any other edge $\{u,v\} \neq \{i,j\}$ in $G$:
\begin{itemize}
    \item If $b_u + b_v = 2$, then $c_u + c_v \ge (r) + (t + 2 - r) = t + 2$, so $a_u + a_v = (c_u + c_v) - 2 \ge t$.
    \item If $b_u + b_v = 1$, then since $c_u + c_v \ge t + 1$, we have $a_u + a_v = (c_u + c_v) - 1 \ge t$.
\end{itemize}
Thus, $\mathbf{a}$ is a $P$-almost $t$-cover of $G$, completing the proof.
\end{proof}

\end{document}
